\begin{proof} Acording to Purpose~\ref{reduction},
we have to show ``$\mbfh_{k-1}\Rightarrow \mbfh_{k}$, for $k=1,\ldots
,n-2$''. So fix an integer $k$, with $1\leq k\leq n-2$ and suppose
that $\mbfh_{k-1}$ holds. Let us prove that $\mbfh_k$ holds. By
Proposition~\ref{gbdesiredform}, we know that there is a Gr\"obner
basis $\tilde{G}_k$, say, of $\ker(\dif_{k})$ of the form: 
\begin{eqnarray}\label{groebf}
\tilde{G}_k:=\bigcup_{i=2}^{r_{k}}\{ \tau_{k}(e_{k,i},e_{k,j})\mid
e_{k,j}\in\mcb_{k,i}\}.
\end{eqnarray}
This is a Gr\"obner basis of $\ker(\dif_{k})$ w.r.t. the monomial
ordering on $\mcc_k$ induced by the monomial ordering on $\mcc_{k-1}$
and the Gr\"obner basis $G_{k-1}$ (see also
Remark~\ref{mon-ordering}). Note that Proposition~\ref{gbdesiredform}
ensures that the cardinality of $\tilde{G}_k$ is $r_{k+1}$.

In Proposition~\ref{tau}, we have shown that, for any $(I_1,\ldots
,I_{k+2})\in\mcb_{k+1}$,
\begin{eqnarray*}
\dif_{k+1}(I_1,\ldots,I_{k+2})=-\tau_{k}(e_{k,i},e_{k,j}),
\end{eqnarray*}
where $e_{k,i}=(I_1,\ldots ,I_{k+1}\cup I_{k+2})$ and
$e_{k,j}=(I_1,\ldots,I_{k}\cup I_{k+1},I_{k+2})$. That is,
$-\dif_{k+1}(I_1,\ldots,I_{k+2})$ is a $\tau$-syzygy associated to the
$S$-vector of $f_{k-1,i}=\dif_{k}(e_{k,i})$ and
$f_{k-1,j}=\dif_{k}(e_{k,j})$. In other words, any element of
$-G_{k}=-\dif_{k+1}(\mcb_{k+1})$ is an element of
$\tilde{G}_k$. Observe that, by Corollary~\ref{ltbij}, $(b)$, the
cardinality of $G_k$ is also $r_{k+1}$. Since the set $-G_{k}$ is
included in the Gr\"obner basis $\tilde{G}_k$ of $\ker(\dif_k)$, and
both $G_k$ and $\tilde{G}_k$ have the same cardinality, it follows
that $G_k=\dif_{k+1}(\mcb_{k+1})$ is a Gr\"obner basis of
$\ker(\dif_{k})$.

(To avoid the preceding discussion using cardinalities, we can argue
as follows. Take $\tau_{k}(e_{k,i},e_{k,j})$, any element in
$\tilde{G}_k$: that is, a $\tau$-syzygy associated to the $S$-vector
of $f_{k-1,i}=\dif_k(e_{k,i})$ and $f_{k-1,j}=\dif_{k}(e_{k,j})$, with
$e_{k,j}\in\mcb_{k,i}$. Remark that, by Discussion~\ref{disc}, both
basis elements $e_{k,i}$ and $e_{k,j}$ are uniquely determined by the
syzygy $\tau_k(e_{k,i},e_{k,j})$. Write $e_{k,i}=(I_1,\ldots
,I_{k},I_{k+1})$ and $e_{k,j}=(I_1,\ldots ,I_{k-1},J_{k},J_{k+1})$,
with $J_{k}\supsetneq I_{k}$. Consider the map
$\varrho_{k,i}:\mcb_{k,i}\to \mcb_{k+1}$ introduced in
Remark~\ref{bki}. Take the element
\begin{eqnarray*}
\varrho_{k,i}(e_{k,j})=(I_1,\ldots,I_{k},J_{k}\setminus
I_{k},J_{k+1})=:(H_1,\ldots,H_k,H_{k+1},H_{k+2})\in\mcb_{k+1}.
\end{eqnarray*}
Observe that $H_s=I_s$, for all $s=1,\ldots ,k$, that
$H_{k+1}=J_{k}\setminus I_{k}$ and that $H_{k+2}=J_{k+1}$. Thus
$H_{k+1}\cup H_{k+2}=(J_{k}\setminus I_{k})\cup J_{k+1}=I_{k+1}$ and
$H_{k}\cup H_{k+1}=I_{k}\cup (J_{k}\setminus I_{k})=J_{k}$. Hence
\begin{eqnarray*}
&& (H_1,\ldots ,H_{k},H_{k+1}\cup H_{k+2})=(I_1,\ldots
  ,I_{k},I_{k+1})=e_{k,i}\mbox{ and }\\ && (H_1,\ldots
  ,H_{k-1},H_{k}\cup H_{k+1},H_{k+2})=(I_1,\ldots
  ,I_{k-1},J_{k},J_{k+1})=e_{k,j}.
\end{eqnarray*}
Then, by Proposition~\ref{tau} again,
$\dif_{k+1}(\varrho_{k,i}(e_{k,j}))=-\tau_{k}(e_{k,i},e_{k,j})$. That
is, any element of $\tilde{G}_k$ is an element of
$-G_{k}=-\dif_{k+1}(\mcb_{k+1})$. In conclusion,
$G_{k}=\dif_{k+1}(\mcb_{k+1})$ is a Gr\"obner basis of
$\ker(\dif_{k})$.)

Since $\mbfh_{k-1}$ holds, it follows that part $(a)$ of
$\mbfh_k$ also holds (recall Notation~\ref{hyphk}). Part $(b)$ of
$\mbfh_k$ follows from Corollary~\ref{ltbij}, $(a)$, and the
hypothesis that $\mbfh_{k-1}$ holds. This completes the proof.
\end{proof}
